% ch2.tex -- Module 8 ESSENTIALS edition, Chapter 2: Structure of the Atom
% (electrons, the nucleus and neutrons). Dot points M8.IQ2.01, M8.IQ2.02
% (NSW Physics Stage 6 Syllabus 2017).
% Knowledge points only: no worked examples, no questions, no answers, no exam tips.
% Condensed from the final full chapter src/mod8/ch2.tex (same notation, same
% facts). TikZ copied from it (captions shortened, a few labels moved to fit).
% The five key experiments use the compact run-in experiment box of the other chapters
% (Aim / What was done / Result / Conclusion and significance), condensed from the full
% book's experiment boxes (validity and reliability lists left out). Each box is short and
% [unbreakable], so no box splits across two pages.
% The Millikan force diagram was moved under the microscope (boxed inset) and the
% Geiger--Marsden figure set beside its text; the physics of both is unchanged.
%
% Chapter macro (prefix chTwo): \chTwoNuc{A}{Z}{X} nuclide symbol, the same as in the
% full book's ch2.
\newcommand\chTwoNuc[3]{\ensuremath{\mathord{%
  \setbox0=\hbox{$\scriptstyle#1$}\setbox2=\hbox{$\scriptstyle#2$}%
  \ifdim\wd0<\wd2 \setbox0=\hbox to\wd2{\hss\box0}\else\setbox2=\hbox to\wd0{\hss\box2}\fi
  {}^{\box0}_{\box2}\mathrm{#3}}}}

\chapter[Structure of the Atom]{Structure of the Atom: Electrons, the Nucleus and Neutrons}\label{ch2:chap}
\covers{M8.IQ2.01, M8.IQ2.02}

\begin{summarybox}{In this chapter}
\begin{itemize}
  \item The electron: cathode rays, Thomson's $q/m$ (1897) and Millikan's charge $e$.
  \item The nucleus: the Geiger--Marsden gold foil and Rutherford's nuclear model (1911).
  \item The neutron: Chadwick's use of conservation laws (1932).
  \item Each step changed the \term{model} of the atom: a simple picture that explains results.
\end{itemize}
\end{summarybox}

%% ==================================================================
\section{Discovering the electron: cathode rays and Thomson}\label{ch2:sec:electron}
\covers{M8.IQ2.01}

\begin{bigidea}
Fields bend cathode rays, so they are negative particles. Thomson's $q/m$ showed that they
are far lighter than atoms: electrons.
\end{bigidea}

\begin{keywords}
  \kw{Cathode ray}{the invisible ray from the cathode ($-$) towards the anode ($+$) in a sealed
    glass tube with almost no air.}
  \kw{Charge-to-mass ratio ($q/m$)}{a particle's charge divided by its mass. Unit: \si{C.kg^{-1}}.}
\end{keywords}

\begin{itemize}
  \item A \textbf{high voltage} (DC) pulls charged particles out of the cathode and speeds them up.
        A \textbf{very low pressure} means few gas particles are in the way.
  \item The glass glows where the rays hit: \term{fluorescence} (giving out light as soon as energy
        hits it). The debate: waves (mostly in Germany) or particles (mostly in Britain).
  \item \textbf{Fits both models:} fluorescence; the sharp shadow of a \textbf{Maltese cross}
        (straight lines). \textbf{Waves:} the rays pass through thin metal foil (Hertz, 1892).
  \item \textbf{Particles:} a \textbf{paddle wheel} in the beam rolls away from the cathode.
        This was taken as evidence that the rays carry momentum. (Today we know the beam
        mainly turns it by heating the blades.) The rays are far slower than light (Thomson, 1894).
  \item \textbf{Negative particles:} a \textbf{magnet} bends the beam like negative charges. A
        metal cup in the beam becomes negative (Perrin, 1895). \textbf{Charged plates} bend it
        towards $+$ (Thomson, 1897). Hertz's plates (1883) did not bend it: gas left in his tube
        became charged and cancelled the field.
  \item \textbf{Decisive:} fields bend charged particles, but not EM (electromagnetic) waves such
        as light (no charge). So cathode rays are streams of negatively charged particles.
\end{itemize}

\subsection*{Thomson's charge-to-mass experiment (1897)}
Anodes with small holes (\term{collimators}) make a thin, straight beam. It passes through
\term{crossed fields}: $E$ and $B$ at \ang{90} to each other and to the beam (\cref{ch2:fig:thomson}).

\begin{figure}
  \centering
  \begin{tikzpicture}[x=1cm, y=1cm, lab/.style={font=\footnotesize, align=center, inner sep=1pt}]
    \draw[line width=0.9pt, hscink, fill=hscvel!5, rounded corners=6pt]
      (3.0,0.42) -- (0,0.42) -- (0,-0.42) -- (3.0,-0.42) -- (3.0,-1.3)
      -- (10.6,-1.3) -- (10.6,1.3) -- (3.0,1.3) -- cycle;
    \draw[hscseriesC, line width=2.6pt] (10.6,-1.15) -- (10.6,1.15);
    \fill[hscvel] (0.3,-0.28) rectangle (0.45,0.28);
    \foreach \x in {1.5,2.3} {
      \fill[hscmuted] (\x,0.07) rectangle (\x+0.12,0.38);
      \fill[hscmuted] (\x,-0.38) rectangle (\x+0.12,-0.07); }
    \draw[hscfield!55, dashed, line width=0.8pt] (6.2,0) circle (1.12);
    \foreach \x in {5.4,6.2,7.0} \foreach \y in {-0.55,0.55}
      \node[text=hscfield, font=\small] at (\x,\y) {$\times$};
    \fill[hscforce] (4.9,0.95) rectangle (7.5,1.08);
    \fill[hscvel] (4.9,-1.08) rectangle (7.5,-0.95);
    \draw[hscvel, line width=1.1pt] (0.5,0) -- (10.55,0);
    \draw[hscvel, line width=0.9pt, dashed] (4.9,0) .. controls (5.77,0) and (6.63,0.083) .. (7.5,0.25) -- (10.55,0.84);
    \draw[hscvel, line width=0.9pt, dashed] (4.9,0) .. controls (5.77,0) and (6.63,-0.083) .. (7.5,-0.25) -- (10.55,-0.84);
    \foreach \y in {-0.84,0,0.84} \fill[hscseriesC!80!black] (10.6,\y) circle (2.4pt);
    \node[lab, anchor=west] at (10.75,0.84) {$E$ only};
    \node[lab, anchor=west] at (10.75,0) {$E$ and $B$\\balanced};
    \node[lab, anchor=west] at (10.75,-0.84) {$B$ only};
    \node[lab, anchor=south] at (0.38,0.5) {cathode};
    \draw[hscmuted, line width=0.4pt] (1.56,0.38) -- (1.8,0.75);
    \draw[hscmuted, line width=0.4pt] (2.36,0.38) -- (2.1,0.75);
    \node[lab, anchor=south] at (1.95,0.72) {anodes\\(collimators)};
    \node[lab, anchor=south, text=hscforce] at (6.2,1.32) {$+$ plate};
    \node[lab, anchor=north, text=hscvel] at (6.2,-1.32) {$-$ plate};
    \draw[hscmuted, line width=0.4pt] (7.0,-0.75) -- (8.3,-1.55);
    \node[lab, anchor=north west, text=hscfield] at (8.0,-1.5) {coils: $B$ into the page};
    \node[lab, anchor=south] at (10.6,1.32) {fluorescent screen};
    \draw[line width=0.6pt] (0.38,-0.28) -- (0.38,-1.1) -- (0.85,-1.1);
    \draw[line width=0.6pt] (2.36,-0.38) -- (2.36,-1.1) -- (1.45,-1.1);
    \draw[line width=0.6pt] (1.56,-0.38) -- (1.56,-1.1);
    \draw[object] (0.85,-1.27) rectangle (1.45,-0.93);
    \node[lab, anchor=north] at (1.15,-1.32) {high voltage};
    \node[font=\small\bfseries, text=hscvel] at (0.7,-0.88) {$-$};
    \node[font=\small\bfseries, text=hscforce] at (1.98,-0.88) {$+$};
  \end{tikzpicture}
  \caption{Thomson's apparatus: crossed fields, then $B$ alone.}\label{ch2:fig:thomson}
\end{figure}

\begin{itemize}
  \item \textbf{Directions:} the electric force is up, towards $+$. With $B$ into the page, the
        right-hand rule pushes a \emph{positive} charge up, so electrons are pushed down.
  \item \textbf{Step 1:} adjust the fields until the beam goes straight: $qE = qvB$, so $v = E/B$.
        Only one speed goes straight: a \term{velocity selector}.
  \item \textbf{Step 2:} switch off $E$. The magnetic force is at \ang{90} to $v$, so the beam
        follows an arc of radius $r$: $qvB = mv^2/r$. (The weight is over $10^{14}$ times smaller
        than $qE$: ignore it.)
  \item \textbf{School e/m tube:} an \term{electron gun} (a hot cathode and an anode) speeds
        electrons from rest through $V$
        ($qV = \tfrac12 mv^2$). A uniform $B$ bends the beam into a full circle; a little gas glows
        along it.
\end{itemize}

\textbf{Why it mattered:} $q/m$ was about 1800 times the hydrogen ion's: for the same size of
charge, about $\tfrac{1}{1800}$ of a hydrogen atom's mass (a \term{subatomic particle}). It was the
same for every metal and gas, so this \term{electron} is in all atoms. Atoms are neutral, so they
must also hold positive charge.

\begin{formulacard}[notonsheet]{Thomson's equations and the school e/m tube}
  \cardformula{v = \frac{E}{B} \qquad\quad \frac{q}{m} = \frac{E}{B^2 r} \qquad\quad
    \frac{q}{m} = \frac{2V}{B^2 r^2}}
  \begin{symbols}
    \sym{$v$}{speed of the particles that go straight}{\si{m.s^{-1}}}
    \sym{$E$, $B$}{electric and magnetic field strengths}{\si{V.m^{-1}}, \si{T}}
    \sym{$q/m$}{charge-to-mass ratio}{\si{C.kg^{-1}}}
    \sym{$r$, $V$}{radius of the beam's path; voltage that speeds it up}{\si{m}, \si{V}}
  \end{symbols}
  \whentouse{crossed fields (first two); sped up from rest through $V$ (third).}
  \derive{$qvB = qE$ gives $v = E/B$. $qvB = mv^2/r$ gives $q/m = v/(Br)$; put in $v = E/B$.
    Tube: $qV = \tfrac12 mv^2$ and $v = (q/m)Br$; remove $v$.}
  \trap{$E$ over $B$, not $B$ over $E$. $r$ in metres; if a scale shows the diameter, halve it.}
\end{formulacard}

%% ==================================================================
\section{Millikan's oil drop experiment}\label{ch2:sec:millikan}
\covers{M8.IQ2.01}

\begin{bigidea}
Every oil drop's charge was a whole number times \qty{1.6e-19}{C}. So charge comes in packets:
the charge of one electron.
\end{bigidea}

\begin{keywords}
  \kw{Quantised}{can only have certain separate values, with nothing in between.}
  \kw{Elementary charge ($e$)}{the smallest charge found on its own, \qty{1.602e-19}{C}: the size
    of an electron's charge.}
\end{keywords}

% Millikan figure set as a full-width figure (not side by side), so it fits under the key
% words and the bullets below it run the full width
\begin{figure}
  \centering
  \begin{tikzpicture}[x=1cm, y=1cm, lab/.style={font=\footnotesize, align=center, inner sep=1pt}]
    \draw[hscink, line width=0.6pt, fill=hscmuted!10] (1.6,2.25) rectangle (5.4,3.6);
    \fill[hscforce!80] (0.55,3.05) circle (0.38);
    \draw[hscink, line width=0.6pt, fill=hscmuted!25] (0.9,2.98) rectangle (1.6,3.12);
    \foreach \p in {(2.1,3.0),(2.6,3.3),(3.0,2.75),(3.6,3.15),(4.2,2.85),(4.7,3.35),(3.3,2.5),(4.9,2.6)}
      \fill[hscacc!80] \p circle (1.6pt);
    \node[lab, anchor=south] at (0.55,3.45) {atomiser};
    \node[lab, anchor=south] at (3.5,3.62) {oil drops};
    \fill[hscforce] (1.4,2.08) rectangle (3.35,2.22);
    \fill[hscforce] (3.65,2.08) rectangle (5.6,2.22);
    \fill[hscvel] (1.4,0.58) rectangle (5.6,0.72);
    \node[font=\small\bfseries, text=hscforce] at (1.25,2.42) {$+$};
    \node[font=\small\bfseries, text=hscvel] at (1.25,0.42) {$-$};
    \foreach \x in {1.75,2.35,4.6,5.1} \draw[field] (\x,2.03) -- (\x,0.77);
    \node[lab, text=hscfield, anchor=east] at (5.06,0.95) {$E$};
    \fill[hscacc] (3.5,1.25) circle (2.6pt);
    \node[lab, anchor=south] at (3.5,1.45) {charged drop};
    \draw[<->, hscink, line width=0.6pt] (5.42,0.72) -- (5.42,2.08) node[pos=0.45, right, inner sep=2pt] {$d$};
    \draw[object] (6.15,1.98) rectangle (8.2,2.42);
    \node[font=\footnotesize] at (7.175,2.2) {X-ray source};
    \draw[hscseriesB, line width=0.9pt, -{Stealth[length=1.8mm]}] (6.13,2.1) -- (5.62,1.8);
    \draw[hscink, line width=0.6pt, fill=hscacc!25] (-0.25,1.2) rectangle (0.25,1.5);
    \node[lab, anchor=north] at (0.0,1.15) {light};
    \draw[hscacc!70, line width=0.6pt, -{Stealth[length=1.6mm]}] (0.3,1.35) -- (1.3,1.35);
    \draw[hscink, line width=0.6pt, fill=hscmuted!20] (6.4,1.22) rectangle (7.9,1.5);
    \draw[hscink, line width=0.6pt, fill=hscmuted!35] (6.25,1.27) rectangle (6.4,1.45);
    \node[lab, anchor=south] at (7.15,1.52) {microscope};
    \draw[line width=0.6pt] (1.4,2.15) -- (1.0,2.15) -- (1.0,-0.2) -- (2.6,-0.2);
    \draw[line width=0.6pt] (3.5,0.58) -- (3.5,0.0);
    \draw[object] (2.6,-0.45) rectangle (4.4,0.0);
    \node[font=\footnotesize] at (3.5,-0.23) {adjustable $V$};
    % inset: forces on the still drop (moved under the microscope so the figure fits beside
    % the text; boxed so it is not read as a second drop)
    \draw[hscmuted, line width=0.4pt, rounded corners=2pt] (6.0,-0.9) rectangle (8.4,1.05);
    \node[font=\scriptsize, text=hscmuted, anchor=north west, inner sep=2pt] at (6.0,1.05) {forces on the drop};
    \begin{scope}[xshift=6.45cm, yshift=-0.1cm]
      \draw[force] (0,0.1) -- (0,0.5) node[right, text=hscforce, inner sep=2pt] {$F_\mathrm{E} = qE$};
      \draw[force] (0,-0.1) -- (0,-0.5) node[right, text=hscforce, inner sep=2pt] {$F_\mathrm{g} = mg$};
      \fill[hscacc] (0,0) circle (3pt);
    \end{scope}
  \end{tikzpicture}
  \caption{Millikan's apparatus and the forces on a still drop. The field points down,
    so it pushes a negative drop up.}\label{ch2:fig:millikan}
\end{figure}

\begin{itemize}
  \item A drop feels its weight $F_\mathrm{g} = mg$ down and the electric force $qE$ up
        (\cref{ch2:fig:millikan}). When it hangs still, $qE = mg$.
  \item \textbf{Mass:} with the field off, the drop falls at its \term{terminal velocity}
        (steady speed, when air resistance balances weight). Heavier drops fall faster.
  \item Oil is used because water evaporates.
  \item Drops gain or lose whole electrons, so the common step between charges is $e$ (the
        smallest measured charge may be $2e$). Millikan measured drops, not one electron.
  \item \textbf{Electron mass:} Millikan's $e$ divided by Thomson's $e/m_\mathrm{e}$ gives
        $m_\mathrm{e} = \qty{9.11e-31}{kg}$, about $\tfrac{1}{1836}$ of a proton's mass.
\end{itemize}

\begin{formulacard}[notonsheet]{Millikan's balanced drop and quantised charge}
  \cardformula{q = \frac{mgd}{V} \qquad\qquad q = ne}
  \begin{symbols}
    \sym{$q$, $m$}{charge (size) and mass of the drop}{\si{C}, \si{kg}}
    \sym{$g$}{gravitational acceleration, \qty{9.8}{m.s^{-2}}}{\si{m.s^{-2}}}
    \sym{$d$, $V$}{distance between the plates; voltage across them}{\si{m}, \si{V}}
    \sym{$n$, $e$}{a whole number of electrons; elementary charge}{no unit, \si{C}}
  \end{symbols}
  \whentouse{a still drop between horizontal plates; $q = ne$ to count electrons.}
  \derive{$qE = mg$ and $E = V/d$ (on the sheet). $q = ne$ is Millikan's result.}
  \trap{not for a moving drop. $d$ in metres. Use the size of $e$ (the sheet prints $-e$).}
\end{formulacard}

%% ==================================================================
\section{The nuclear atom: Geiger--Marsden and Rutherford}\label{ch2:sec:nucleus}
\covers{M8.IQ2.02}

\begin{bigidea}
Most alphas went straight through gold foil, but a very few bounced back. So an atom is
mostly empty space, with a tiny, dense, positive nucleus.
\end{bigidea}

\begin{keywords}
  \kw{Plum pudding model}{Thomson's model (1904): spread-out positive charge with electrons
    dotted through it. No nucleus.}
  \kw{Alpha particle ($\alpha$)}{a helium nucleus: 2 protons and 2 neutrons, charge $+2e$.}
\end{keywords}

\Cref{ch2:fig:models} compares the plum pudding model with Rutherford's \term{nuclear model} (1911).

\begin{figure}
  \centering
  \begin{tikzpicture}[x=1cm, y=1cm, lab/.style={font=\footnotesize, align=left, inner sep=1pt}]
    \begin{scope}
      \fill[hscforce!13] (0,0) circle (1.3);
      \draw[hscforce!60, line width=0.8pt] (0,0) circle (1.3);
      \foreach \p in {(-0.6,0.55),(0.15,0.8),(0.7,0.35),(-0.85,-0.2),(-0.1,0.05),(0.55,-0.5),(-0.35,-0.75)}
        \node[text=hscforce, font=\small\bfseries] at \p {$+$};
      \foreach \p in {(-0.25,0.5),(0.45,0.0),(-0.6,-0.5),(0.2,-0.85),(0.8,0.75),(-0.95,0.25),(0.15,-0.35)} {
        \fill[hscvel] \p circle (3.2pt);
        \node[text=white, font=\tiny] at \p {$-$}; }
      \draw[hscmuted, line width=0.4pt] (0.98,0.52) -- (1.75,1.05);
      \node[lab, anchor=west] at (1.75,1.05) {positive charge spread\\through the whole atom};
      \draw[hscmuted, line width=0.4pt] (0.48,0.03) -- (1.75,-0.35);
      \node[lab, anchor=west] at (1.75,-0.35) {electrons};
      \node[font=\footnotesize, anchor=north] at (1.2,-1.4) {(a) Plum pudding model (Thomson, 1904)};
    \end{scope}
    \begin{scope}[xshift=8.2cm]
      \draw[hscmuted, dashed, line width=0.6pt] (0,0) circle (1.3);
      \fill[hscforce] (0,0) circle (1.6pt);
      \foreach \p in {(-0.7,0.6),(0.55,0.9),(0.3,-0.95),(-0.95,-0.45),(0.95,-0.25),(-0.25,-0.8)} {
        \fill[hscvel] \p circle (3.2pt);
        \node[text=white, font=\tiny] at \p {$-$}; }
      \draw[hscmuted, line width=0.4pt] (0.05,0.03) -- (1.75,0.95);
      \node[lab, anchor=west] at (1.75,0.95) {tiny nucleus ($+$)};
      \node[lab, anchor=west] at (1.75,0.2) {mostly empty space};
      \draw[hscmuted, line width=0.4pt] (0.3,-0.95) -- (1.75,-0.55);
      \node[lab, anchor=west] at (1.75,-0.55) {electrons around\\the nucleus};
      \node[font=\footnotesize, anchor=north] at (1.2,-1.4) {(b) Nuclear model (Rutherford, 1911)};
    \end{scope}
  \end{tikzpicture}
  \caption{Two models of the atom. In (b) the nucleus is drawn far too big.}\label{ch2:fig:models}
\end{figure}

\begin{sidebyside}[0.6]
  \centering
  \begin{tikzpicture}[x=1cm, y=1cm, lab/.style={font=\footnotesize, align=left, inner sep=1pt}]
    \draw[hscmuted, dashed, line width=0.6pt] (0,0) circle (2.55);
    \draw[hscseriesC, line width=2.2pt] (-168:2.05) arc (-168:168:2.05);
    \fill[hscmuted!45] (-4.3,-0.5) rectangle (-3.5,0.5);
    \fill[white] (-3.85,-0.07) rectangle (-3.5,0.07);
    \fill[hscacc] (-3.95,0) circle (3pt);
    \draw[hscacc, line width=1.1pt, -{Stealth[length=2mm]}] (-3.45,0) -- (-0.08,0);
    \draw[draw=yellow!60!black, line width=2.4pt] (0,-0.55) -- (0,0.55);
    \foreach \a in {-3,0,3} \draw[hscacc, line width=0.9pt, -{Stealth[length=1.8mm]}] (0,0) -- (\a:1.95);
    \draw[hscacc, line width=0.9pt, -{Stealth[length=1.8mm]}] (0,0) -- (25:1.95);
    \draw[hscacc, line width=0.9pt, -{Stealth[length=1.8mm]}] (0,0) -- (-30:1.95);
    \draw[hscacc, line width=0.9pt, -{Stealth[length=1.8mm]}] (0,0) -- (140:1.95);
    \begin{scope}[rotate=55]
      \draw[hscink, line width=0.6pt, fill=hscmuted!20] (2.12,-0.14) rectangle (3.1,0.14);
    \end{scope}
    \node[lab, anchor=north, align=center] at (-3.9,-0.55) {alpha source\\in lead};
    \draw[hscmuted, line width=0.4pt] (-0.04,-0.52) -- (-0.4,-0.74);
    \node[lab, anchor=north] at (-0.6,-0.74) {thin gold foil};
    \node[lab, anchor=west] at (2.7,0.15) {most:\\straight};
    \node[lab, anchor=west] at (2.55,1.45) {some:\\small};
    \draw[hscmuted, line width=0.4pt] (140:1.75) -- (-2.25,2.0);
    \node[lab, anchor=south east] at (-2.2,1.95) {about 1 in 8000:\\more than \ang{90}};
    \node[lab, anchor=south] at (1.6,2.75) {microscope (moves around)};
    \draw[hscmuted, line width=0.4pt] (-45:2.05) -- (2.0,-2.3);
    \node[lab, anchor=west] at (2.0,-2.3) {zinc sulfide\\screen};
    \node[lab, anchor=north] at (-2.2,-2.45) {vacuum chamber};
    \draw[hscmuted, line width=0.4pt] (-2.2,-2.45) -- (-130:2.55);
  \end{tikzpicture}
  \captionof{figure}{The Geiger--Marsden experiment (top view).}\label{ch2:fig:goldfoil}
\nextside
  \textbf{The test (1909):} Geiger and Marsden fired alphas at gold foil about 1000 atoms thick,
  in a vacuum (\cref{ch2:fig:goldfoil}). A microscope counted flashes (\term{scintillations}) on
  a zinc sulfide screen at each angle.

  \textbf{The prediction} (plum pudding): spread-out charge pushes weakly, and electrons are about
  7300 times lighter than an alpha. So all alphas should pass straight through.
\end{sidebyside}

\Cref{ch2:tab:goldfoil} shows what really happened.

\begin{table}
  \caption{The gold-foil results, and how Rutherford explained them (1911).}\label{ch2:tab:goldfoil}
  \centering\small
  \begin{tabular}{@{}L{50mm}L{94mm}@{}}
    \toprule
    \colhead{Observation} & \colhead{Rutherford's explanation}\\
    \midrule
    Most alphas went straight through. & The atom is mostly empty space.\\
    Some were deflected by small angles. & They passed fairly close to a positive nucleus.\\
    About 1 in 8000 were deflected by more than \ang{90}. & They headed almost straight at a
      tiny nucleus with nearly all the mass, and were repelled back. Rare: the nucleus is tiny.\\
    \bottomrule
  \end{tabular}
\end{table}

\begin{itemize}
  \item \textbf{Rutherford's nuclear model (1911):} a tiny nucleus holds all the positive
        charge and more than 99.9\% of the mass. It is about \num{10000} to \num{100000} times smaller
        across than the atom. Electrons move around it, far away. The rest is empty space.
  \item The alphas never touch the nucleus: the electric force repels them. By Coulomb's law
        ($F = \frac{1}{4\pi\varepsilon_0}\frac{q_1 q_2}{r^2}$), only tightly packed charge can turn
        a fast alpha back.
  \item \textbf{Strengths:} it explains all three results and why atoms are neutral. In 1913
        counts at many angles, metals and foil thicknesses matched Coulomb's-law predictions.
  \item \textbf{Limitations:} an orbiting electron accelerates, so classical physics says it should
        radiate, lose energy and spiral in. It cannot explain line spectra. It does not say what
        the nucleus is made of, or why nuclei have about twice the mass of their protons.
  \item \textbf{Judgement:} strongly supported for the nucleus, but incomplete for the electrons.
  \item \textbf{Why this set-up:} gold can be made very thin and has a large nuclear charge
        ($+79e$). The vacuum stops air from scattering the alphas.
\end{itemize}

%% ==================================================================
\section{Chadwick and the neutron}\label{ch2:sec:neutron}
\covers{M8.IQ2.02}

\begin{bigidea}
Nuclei have about twice the mass of their protons. Chadwick (1932) used conservation of
energy and momentum to show that the missing mass is neutrons.
\end{bigidea}

\begin{keywords}
  \kw{Neutron ($\chTwoNuc{1}{0}{n}$)}{a neutral particle in the nucleus with about the mass of a
    proton.}
  \kw{Elastic collision}{a collision in which the total kinetic energy stays the same.}
\end{keywords}

\begin{itemize}
  \item \textbf{Missing mass:} in 1919 Rutherford knocked hydrogen nuclei out of nitrogen with
        alphas. In 1920 he named them \term{protons} (charge $+e$). Helium has 2 protons but
        about 4 times hydrogen's mass. So he predicted a neutral particle (1920).
  \item \textbf{1930:} alphas on beryllium gave a very penetrating radiation that fields did not bend
        (no charge). Bothe and Becker thought it was gamma rays.
  \item \textbf{1932:} the radiation knocks fast protons out of paraffin wax (Curie and Joliot).
        Chadwick also used nitrogen targets (\cref{ch2:fig:chadwick}). An \term{ionisation chamber}
        measured the recoils: a charged particle ionises its gas, making a tiny current.
\end{itemize}

\begin{figure}
  \centering
  \begin{tikzpicture}[x=1cm, y=1cm,
      bx/.style={draw=hscink, line width=0.7pt, fill=hscmuted!12, rounded corners=2pt,
        minimum width=18mm, minimum height=10mm, align=center, font=\footnotesize},
      lab/.style={font=\footnotesize, align=center, inner sep=1pt}]
    \node[bx] (po) at (1.1,0) {polonium\\(alpha source)};
    \node[bx] (be) at (4.85,0) {beryllium};
    \node[bx] (wax) at (8.75,0) {paraffin wax\\(lots of H)};
    \node[bx] (det) at (13.1,0) {ionisation chamber\\(detector)};
    \draw[hscacc, line width=1.1pt, -{Stealth[length=2mm]}] (po.east) -- (be.west)
      node[midway, above=2.5pt, lab] {alpha\\particles};
    \draw[hscmuted, line width=1.1pt, dashed, -{Stealth[length=2mm]}] (be.east) -- (wax.west)
      node[midway, above=2.5pt, lab] {neutral\\radiation};
    \draw[hscforce, line width=1.1pt, -{Stealth[length=2mm]}] (wax.east) -- (det.west)
      node[midway, above=2.5pt, lab] {protons};
    \begin{scope}[yshift=-1.75cm]
      \node[lab, anchor=west] at (-0.05,0.35) {neutron hits a proton (same mass):};
      \fill[hscmuted] (0.3,-0.25) circle (5pt); \node[lab, below=5.5pt] at (0.3,-0.25) {n};
      \draw[velocity] (0.5,-0.25) -- (1.3,-0.25);
      \fill[hscforce] (1.7,-0.25) circle (5pt); \node[lab, below=5.5pt] at (1.7,-0.25) {p};
      \node[font=\small] at (2.45,-0.25) {$\Rightarrow$};
      \fill[hscmuted] (3.2,-0.25) circle (5pt); \node[lab, below=5.5pt] at (3.2,-0.25) {stops};
      \fill[hscforce] (4.1,-0.25) circle (5pt);
      \draw[velocity] (4.3,-0.25) -- (5.1,-0.25);
      \node[lab, anchor=north west] at (-0.05,-0.88) {the proton gets almost all the energy};
    \end{scope}
    \begin{scope}[xshift=7.2cm, yshift=-1.75cm]
      \node[lab, anchor=west] at (-0.05,0.35) {neutron hits nitrogen (14 times the mass):};
      \fill[hscmuted] (0.3,-0.25) circle (5pt); \node[lab, below=5.5pt] at (0.3,-0.25) {n};
      \draw[velocity] (0.5,-0.25) -- (1.3,-0.25);
      \fill[hscseriesA!70] (1.85,-0.25) circle (10pt); \node[font=\scriptsize, text=white] at (1.85,-0.25) {N};
      \node[font=\small] at (2.7,-0.25) {$\Rightarrow$};
      \draw[velocity] (3.75,-0.25) -- (3.05,-0.25);
      \fill[hscmuted] (3.95,-0.25) circle (5pt); \node[lab, below=5.5pt] at (3.95,-0.25) {bounces back};
      \fill[hscseriesA!70] (5.5,-0.25) circle (10pt); \node[font=\scriptsize, text=white] at (5.5,-0.25) {N};
      \draw[velocity] (5.9,-0.25) -- (6.15,-0.25);
      \node[lab, anchor=north west] at (-0.05,-0.88) {nitrogen gets only about a quarter of the energy};
    \end{scope}
  \end{tikzpicture}
  \caption{Chadwick's experiment (top) and head-on elastic collisions (bottom).}\label{ch2:fig:chadwick}
\end{figure}

\begin{itemize}
  \item \textbf{Energy rules out gamma rays:} the fastest protons had about \qty{5.7}{MeV}. A
        photon has no mass and little momentum for its energy, so it would need about
        \qty{50}{MeV}. The beryllium reaction could give at most about \qty{14}{MeV}.
  \item \textbf{Momentum and kinetic energy give the mass} (both on the sheet):
        $\sum m\vec{v}_\mathrm{before} = \sum m\vec{v}_\mathrm{after}$ (no outside force) and
        $\sum\frac{1}{2}mv^2_\mathrm{before} = \sum\frac{1}{2}mv^2_\mathrm{after}$ (elastic only).
        They give a target's recoil speed from the particle's unknown mass $m$ and unknown
        speed $u$. Chadwick used two targets: hydrogen (about 1 atomic mass unit, \qty{1}{u})
        and nitrogen (about \qty{14}{u}). Dividing the two results cancels the unknown speed
        $u$. Only the mass is left: $m \approx 1.15$ atomic mass units, about a proton's mass.
  \item The reaction: $\chTwoNuc{4}{2}{He} + \chTwoNuc{9}{4}{Be} \rightarrow \chTwoNuc{12}{6}{C}
        + \chTwoNuc{1}{0}{n}$. In $\chTwoNuc{A}{Z}{X}$, $A$ = protons + neutrons and $Z$ = protons;
        the totals of $A$ and of $Z$ are the same on both sides. The neutron
        has no charge and leaves no track: this is \term{indirect evidence} (from effects only).
\end{itemize}


%% ==================================================================
\section{Key experiments}\label{ch2:sec:experiments}
\covers{M8.IQ2.01, M8.IQ2.02}

To \textbf{model} an experiment is to draw and label its set-up, or to use a simulation or an
analogy. One analogy: marbles rolled under a board reveal the shape of a hidden object.

\begin{experiment}[unbreakable]{Early cathode-ray experiments (1870s--1890s)}
\textbf{Aim:} to find whether cathode rays are waves or particles. \textbf{What was done:}
objects (a Maltese cross, a paddle wheel) were put in the beam. Charged plates or a magnet were
brought near the tube. \textbf{Result:} sharp shadows; fields bent the beam like negative
charges. \textbf{Conclusion and significance:} the rays travel in straight lines and carry
negative charge: they are negative particles. This led to Thomson's work.
\end{experiment}

\begin{experiment}[unbreakable]{Thomson's charge-to-mass experiment (1897)}
\textbf{Aim:} to measure $q/m$ for cathode-ray particles. \textbf{What was done:} $E$ and $B$
were adjusted until the beam went straight ($v = E/B$). Then $E$ was switched off; the radius
$r$ gave $q/m = E/(B^2 r)$. This was repeated with other metals and gases. \textbf{Result:}
about \qty{e11}{C.kg^{-1}} (today \qty{1.759e11}{C.kg^{-1}}) for every metal and gas.
\textbf{Conclusion and significance:} identical negative particles, far lighter than atoms, are
in all matter: electrons. This first subatomic particle led to the plum pudding model.
\end{experiment}

\begin{experiment}[unbreakable]{Millikan's oil drop experiment (1909--1913)}
\textbf{Aim:} to find whether charge comes in a smallest unit, and measure it. \textbf{What
was done:} with the field off, a drop's falling speed gave its mass. Then $V$ was adjusted
until the drop hung still: $q = mgd/V$. X-rays changed its charge; many drops were
measured. \textbf{Result:} every charge was a whole number times \qty{1.592e-19}{C} (today
\qty{1.602e-19}{C}). \textbf{Conclusion and significance:} charge is quantised; the smallest
unit $e$ is the electron's charge. With Thomson's $q/m$, it gave the electron's mass.
\end{experiment}

\begin{experiment}[unbreakable]{The Geiger--Marsden experiment (1909)}
\textbf{Aim:} to test the plum pudding model. \textbf{What was done:} a narrow beam of alphas
hit thin gold foil in a vacuum. Flashes on a zinc sulfide screen were counted through a
microscope at each angle. \textbf{Result:} most went straight through; some had small
deflections; about 1 in 8000 turned by more than \ang{90}. \textbf{Conclusion and
significance:} a tiny, dense, positive nucleus; the atom is mostly empty space. This was the
discovery of the nucleus, the basis of Bohr's model and of nuclear physics.
\end{experiment}

\begin{experiment}[unbreakable]{Chadwick's discovery of the neutron (1932)}
\textbf{Aim:} to identify the neutral radiation from beryllium hit by alphas. \textbf{What was
done:} the radiation hit paraffin wax, then nitrogen. The recoil energies of the protons and
nitrogen nuclei were measured. Conservation of energy and momentum were applied.
\textbf{Result:} protons up to \qty{5.7}{MeV}; nitrogen nuclei recoiled slowly; no single
gamma-ray energy fits both. \textbf{Conclusion and significance:} neutral particles with about
the proton's mass: neutrons. This completed the picture of the atom and explained the missing
mass. Neutrons were later used to split nuclei.
\end{experiment}

%% ==================================================================
\begin{summarybox}{Remember}
\begin{itemize}
  \item Fields bend cathode rays but not light, so the rays are negative particles.
  \item Thomson: $v = E/B$, then $q/m = E/(B^2 r)$ (e/m tube: $2V/(B^2 r^2)$), the same for all
        metals and gases: electrons are in all atoms. Not on the sheet.
  \item Millikan: a still drop has $qE = mg$, so $q = mgd/V$. Every charge is $q = ne$: quantised.
  \item Gold foil: most alphas pass straight through (empty space); about 1 in 8000 turn by more
        than \ang{90} (a tiny, dense, positive nucleus).
  \item Rutherford's model is strong for the nucleus, but it cannot explain line spectra or why
        orbiting electrons do not spiral in.
  \item Chadwick: a gamma ray would need more energy than the reaction gives. Momentum and
        kinetic energy, with two targets, give a neutral particle with a proton's mass.
\end{itemize}
\end{summarybox}
